% from ega0-20-fr.tex
\oldpage[0\textsubscript{IV-1}]{145}
d'où, par passage à l'homologie, un \emph{diagramme commutatif}
\begin{equation}
\tag{20.6.21.3}
\label{0.20.6.21.3-fr}
\xymatrix@C=0.7em{
  0 \ar[r] & \Upsilon^C_{B/A/\Lambda} \ar[r] \ar[d] &
  \Upsilon_{C/A/\Lambda} \ar[r] \ar[d] &
  \Upsilon_{C/B/\Lambda} \ar[r] \ar[d] &
  \Omega^1_{B/A}\otimes_B C \ar[r] \ar[d] &
  \Omega^1_{C/A} \ar[r] \ar[d] &
  \Omega^1_{C/B} \ar[r] \ar[d] & 0 \\
  0 \ar[r] & \Upsilon^{C'}_{B'/A'/\Lambda'} \ar[r] &
  \Upsilon_{C'/A'/\Lambda'} \ar[r] &
  \Upsilon_{C'/B'/\Lambda'} \ar[r] &
  \Omega^1_{B'/A'}\otimes_{B'} C' \ar[r] &
  \Omega^1_{C'/A'} \ar[r] & \Omega^1_{C'/B'} \ar[r] & 0
}
\end{equation}
On a un diagramme commutatif analogue pour \sref{0.20.6.18.1-fr}.
\end{enumerate}
\end{remarks}

\begin{proposition}[20.6.22]
\label{0.20.6.22-fr}
Soient $s:\Lambda\to A$, $u:A\to B$ deux homomorphismes d'anneaux,
$\mathfrak{R}$ un idéal de $B$, $C$ l'anneau quotient $B/\mathfrak{R}$.
Supposons que $E=B/\mathfrak{R}^2$ soit une $B$-extension
$\Lambda$-triviale de $C$ par $\mathfrak{R}/\mathfrak{R}^2$. On a alors une
suite exacte
\begin{equation}
\tag{20.6.22.1}
\label{0.20.6.22.1-fr}
  0\to\Upsilon^C_{B/A/\Lambda}\xrightarrow{v'}
  \Upsilon_{C/A/\Lambda}\xrightarrow{\chi_E}
  \mathfrak{R}/\mathfrak{R}^2\xrightarrow{\delta_{C/B/A}}
  \Omega^1_{B/A}\otimes_B C\xrightarrow{v_{C/B/A}}
  \Omega^1_{C/A}\to0.
\end{equation}
\end{proposition}

\begin{proof}
En effet, comme $v:B\to C$ est surjectif, on a $\Omega^1_{C/B}=0$
\sref{0.20.4.12-fr}. En outre, il résulte de \sref{0.20.6.10-fr} que
$\Upsilon_{C/B/\Lambda}$ s'identifie canoniquement à
$\mathfrak{R}/\mathfrak{R}^2$. Il suffit alors d'appliquer les suites exactes
\sref{0.20.6.17.1-fr} et \sref{0.20.6.18.1-fr}, en notant que l'on a un
diagramme commutatif
\[
\xymatrix{
  \Upsilon_{C/A/\Lambda} \ar[r] \ar[d]_{\chi_E} &
  \Omega^1_{A/\Lambda}\otimes_A C \ar[d] \\
  \Upsilon_{C/E/\Lambda}=\Upsilon_{C/B/\Lambda}=\mathfrak{R}/\mathfrak{R}^2
    \ar[r]^-{\delta_{C/B/A}} \ar[d]_{\ell} &
  \Omega^1_{B/A}\otimes_B C=\Omega^1_{E/A}\otimes_E C \ar[d] \\
  \mathfrak{R}/\mathfrak{R}^2 \ar[r]^{\delta_{C/B/A}} &
  \Omega^1_{B/A}\otimes_B C
}
\]
et utilisant la commutativité du diagramme \sref{0.20.6.17.3-fr}.
\end{proof}

On a ainsi précisé le \emph{noyau} de $\delta_{C/B/A}$ dans le cas où il
existe un anneau $\Lambda$ tel que $A$ soit une $\Lambda$-algèbre et que
$B/\mathfrak{R}^2$ soit $\Lambda$-triviale ; on observera que c'est le cas
en particulier lorsque $C$ est une \emph{$\Lambda$-algèbre formellement
lisse} (pour la topologie discrète).

Supposons en outre qu'on ait un diagramme commutatif d'homomorphismes
d'anneaux
\[
\xymatrix{
  \Lambda' \ar[r] & A' \ar[r] & B' \\
  \Lambda \ar[u] \ar[r] & A \ar[u] \ar[r] & B \ar[u]^f
}
\]
